\section{Especificación de Requerimientos del Sistema}

Este capítulo detalla los requerimientos formales del sistema Olympus. Siguiendo los estándares de ingeniería de sistemas (ISO/IEC/IEEE 29148), cada requerimiento utiliza el verbo auxiliar \textbf{``deberá''} para denotar obligatoriedad.

\subsection{Requerimientos de Capacidad (Funcionales)}

\subsubsection{Navegación y Control (GNC)}
\begin{itemize}[label=---]
    \item \textbf{RF-GNC-01:} El sistema \textbf{deberá} calcular rutas de navegación autónoma evitando obstáculos mayores a 10 cm de altura. \\
    \textit{[Prioridad: Crítica | Fuente: Misión | V\&V: Prueba (T)]}
    \item \textbf{RF-GNC-02:} El rover \textbf{deberá} ser capaz de realizar una parada de emergencia de forma autónoma ante la pérdida de datos del sensor LiDAR. \\
    \textit{[Prioridad: Crítica | Fuente: Seguridad | V\&V: Demostración (D)]}
    \item \textbf{RF-GNC-03:} El sistema \textbf{deberá} mantener una precisión de localización con un error máximo de $\pm$ 50 cm en un área de 10x10 metros. \\
    \textit{[Prioridad: Alta | Fuente: Misión | V\&V: Análisis (A)]}
\end{itemize}

\subsubsection{Movilidad}
\begin{itemize}[label=---]
    \item \textbf{RF-MOV-01:} El tren de rodaje \textbf{deberá} permitir el avance en pendientes de hasta 15$^\circ$ sobre regolito simulado. \\
    \textit{[Prioridad: Crítica | Fuente: Diseño | V\&V: Prueba (T)]}
    \item \textbf{RF-MOV-02:} El rover \textbf{deberá} alcanzar una velocidad máxima nominal de 10 cm/s en terreno llano. \\
    \textit{[Prioridad: Media | Fuente: Operaciones | V\&V: Prueba (T)]}
\end{itemize}

\subsubsection{Gestión de Energía}
\begin{itemize}[label=---]
    \item \textbf{RF-ENE-01:} El subsistema de potencia \textbf{deberá} monitorizar el voltaje y la corriente de las baterías con una frecuencia mínima de 1 Hz. \\
    \textit{[Prioridad: Alta | Fuente: Diseño | V\&V: Inspección (I)]}
    \item \textbf{RF-ENE-02:} El sistema \textbf{deberá} iniciar un protocolo de ``Modo Seguro'' cuando el estado de carga (SoC) sea inferior al 15\%. \\
    \textit{[Prioridad: Crítica | Fuente: Seguridad | V\&V: Demostración (D)]}
\end{itemize}

\subsection{Requerimientos de Seguridad y FDIR (Fault Detection, Isolation, and Recovery)}
\begin{itemize}[label=---]
    \item \textbf{RF-FDI-01:} El sistema \textbf{deberá} monitorizar el latido (\textit{heartbeat}) de todos los nodos ROS 2 críticos y reiniciar aquellos que fallen tras 3 intentos. \\
    \textit{[Prioridad: Crítica | Fuente: Seguridad | V\&V: Demostración (D)]}
    \item \textbf{RF-FDI-02:} El rover \textbf{deberá} entrar en estado de reposo (Safe Mode) si se pierde el enlace de mando por más de 10 segundos. \\
    \textit{[Prioridad: Crítica | Fuente: Seguridad | V\&V: Prueba (T)]}
\end{itemize}

\subsection{Interfaces de Sistemas}

\subsubsection{Interfaces de Hardware}
\begin{itemize}[label=---]
    \item \textbf{RI-HW-01:} La unidad central \textbf{deberá} conectarse con el LiDAR a través de una interfaz USB 3.0 o Ethernet. \\
    \textit{[Prioridad: Alta | Fuente: Diseño | V\&V: Inspección (I)]}
    \item \textbf{RI-HW-02:} El bus de motores \textbf{deberá} utilizar la interfaz física CAN con una velocidad mínima de 500 kbps. \\
    \textit{[Prioridad: Crítica | Fuente: Diseño | V\&V: Prueba (T)]}
\end{itemize}

\subsubsection{Interfaces de Software}
\begin{itemize}[label=---]
    \item \textbf{RI-SW-01:} El sistema \textbf{deberá} utilizar el middleware ROS 2 (Humble) para la comunicación entre nodos internos. \\
    \textit{[Prioridad: Crítica | Fuente: Diseño | V\&V: Inspección (I)]}
    \item \textbf{RI-SW-02:} Los mensajes de telemetría \textbf{deberán} encapsularse en datagramas CSP (CubeSat Space Protocol) para su transmisión inalámbrica. \\
    \textit{[Prioridad: Crítica | Fuente: Diseño | V\&V: Demostración (D)]}
\end{itemize}

\subsection{Atributos de Calidad (No Funcionales)}

\subsubsection{Confiabilidad y Rendimiento}
\begin{itemize}[label=---]
    \item \textbf{RNF-PER-01:} La latencia de la telemetría enviada mediante CSP desde el rover a la GCS no \textbf{deberá} exceder los 500 ms. \\
    \textit{[Prioridad: Media | Fuente: Operaciones | V\&V: Prueba (T)]}
    \item \textbf{RNF-PER-02:} El bucle de control Nav2 \textbf{deberá} ejecutarse a una frecuencia mínima de 10 Hz. \\
    \textit{[Prioridad: Alta | Fuente: Diseño | V\&V: Análisis (A)]}
\end{itemize}

\subsubsection{Ciberseguridad}
\begin{itemize}[label=---]
    \item \textbf{RNF-SEC-01:} El enlace de comunicaciones \textbf{deberá} implementar un cifrado AES-128 para prevenir la interceptación de datos. \\
    \textit{[Prioridad: Alta | Fuente: Seguridad | V\&V: Inspección (I)]}
\end{itemize}
